\textbf{1. Kiến thức cần nhớ}
\begin{itemize}
\item \textbf{Bất phương trình bậc nhất một ẩn} là bất phương trình có dạng:
$$\boxed{ax + b < 0\quad (\text{hoặc } ax + b > 0,\ ax + b \le 0,\ ax + b \ge 0)}$$
trong đó $a, b$ là hai số đã cho và $a \ne 0$.
\item \textbf{Quy tắc chuyển vế}: Khi chuyển một hạng tử của bất phương trình từ vế này sang vế kia, ta phải \textbf{đổi dấu} hạng tử đó.
\item \textbf{Quy tắc nhân với một số}:
\begin{itemize}
\item Nhân (hoặc chia) hai vế với một số \textbf{dương} $\rightarrow$ \textbf{giữ nguyên chiều}.
\item Nhân (hoặc chia) hai vế với một số \textbf{âm} $\rightarrow$ \textbf{đổi chiều} bất phương trình.
\end{itemize}
\end{itemize}

\textbf{2. Phương pháp giải các dạng}

%% DANG 9D2?6-1
\textbf{Dạng 1. Nhận biết bất phương trình bậc nhất một ẩn, kiểm tra nghiệm}
\begin{enumerate}
\item Nhận dạng: Ẩn số $x$ có số mũ là $1$ và hệ số $a \ne 0$.
\item Để kiểm tra $x_0$ có là nghiệm không: thay $x = x_0$ vào bất phương trình. Nếu thu được khẳng định đúng thì $x_0$ là nghiệm; ngược lại không phải nghiệm.
\end{enumerate}

%% DANG 9D2?6-2
\textbf{Dạng 2. Giải bất phương trình bậc nhất một ẩn}
\begin{enumerate}
\item Chuyển các hạng tử chứa ẩn sang vế trái, các hằng số sang vế phải: $ax > -b$.
\item Chia cả hai vế cho hệ số $a$:
\begin{itemize}
\item Nếu $a > 0 \implies \boxed{x > -\dfrac{b}{a}}$ (giữ nguyên chiều).
\item Nếu $a < 0 \implies \boxed{x < -\dfrac{b}{a}}$ (\textbf{đổi chiều} bất phương trình).
\end{itemize}
\item Biểu diễn tập nghiệm trên trục số (dùng ngoặc tròn $($ hoặc ngoặc vuông $[$, gạch bỏ phần không thuộc tập nghiệm).
\end{enumerate}

%% DANG 9D2?6-3
\textbf{Dạng 3. Bất phương trình quy về bậc nhất một ẩn}
\begin{enumerate}
\item Bỏ dấu ngoặc (nhân đơn thức, đa thức hoặc dùng hằng đẳng thức).
\item Quy đồng và khử mẫu số dương (nhân hai vế với BCNN của các mẫu số nguyên dương).
\item Thu gọn các hạng tử đồng dạng để đưa về dạng $ax > c$ hoặc $ax < c$.
\item Giải và kết luận tập nghiệm.
\end{enumerate}

%% DANG 9D2?6-4
\textbf{Dạng 4. Bài toán thực tế dẫn đến bất phương trình}
\begin{enumerate}
\item Chọn ẩn số $x$ và đặt điều kiện thích hợp.
\item Biểu diễn đại lượng thực tế theo $x$.
\item Lập bất phương trình từ các điều kiện ràng buộc trong đề bài (ví dụ: ``kinh phí không quá $5$ triệu đồng'', ``điểm trung bình ít nhất $8{,}0$''\dots).
\item Giải bất phương trình và chọn các giá trị của $x$ thỏa mãn điều kiện thực tế (thường là số nguyên dương).
\end{enumerate}
